\chapter{Simulazione e Analisi Pratica}
\label{cap:simulazione}

Questo capitolo descrive la parte operativa del lavoro: come sono stati trattati i dati
pubblici del GME, come sono state ricostruite le curve di domanda e offerta, e con quale
accuratezza la ricostruzione riproduce i prezzi effettivamente formatisi sul mercato. La
validazione \`e presentata prima degli scenari di accumulo perch\'e ne costituisce il
presupposto: un modello che non sa riprodurre il prezzo osservato non pu\`o essere usato per
calcolare un prezzo controfattuale.

\section{Setup della simulazione}
\label{sec:setup}

\subsection{Ricostruzione delle curve aggregate}

Per ogni asta — individuata da una zona e da un periodo — le offerte di vendita sono ordinate
per prezzo crescente e quelle di acquisto per prezzo decrescente, e le rispettive quantit\`a
sono cumulate. Si ottengono cos\`i le due funzioni a gradini $S(p)$ e $D(p)$ introdotte nella
Sezione~\ref{sec:mgp}, il cui punto di incontro definisce il prezzo di
equilibrio~\eqref{eq:clearing}.

Una precisazione sul perch\'e le curve siano costruite dalle \textbf{offerte} e non dal merit
order degli impianti. L'alternativa naturale sarebbe ordinare il parco di generazione per costo
marginale crescente, ricavandolo da combustibile, rendimento e costi della CO\textsubscript{2}.
\textcite{veenstra2025profitability} argomentano per\`o che sul breve periodo la relazione fra il
merit order costruito sugli impianti e le offerte effettivamente presentate al mercato \`e
\textbf{debole}, perch\'e un'offerta non corrisponde necessariamente a uno specifico impianto:
gli operatori presentano portafogli, aggregano capacit\`a in impianti virtuali e rivendono
energia acquistata altrove.

Ne segue che il merit order tecnico descrive il sistema fisico ma non la curva su cui il prezzo
si forma davvero. Poich\'e l'oggetto di questa tesi \`e precisamente \textbf{come l'accumulo
sposta il prezzo di equilibrio}, la rappresentazione pertinente \`e quella delle offerte, che
\`e anche l'unica osservabile senza stimare parametri di costo non pubblici. La scelta ha un
costo — le curve ricostruite non dicono nulla sulla tecnologia marginale — ma \`e quello giusto
da pagare per la domanda posta.

L'implementazione valuta le curve sull'insieme di tutti i prezzi presenti nell'asta.
L'insieme \`e sufficiente perch\'e le curve sono costanti a tratti e cambiano valore solo in
corrispondenza di un prezzo offerto: l'intersezione cade quindi sempre su uno di questi
punti. Il costo computazionale \`e $O(n \log n)$ nel numero $n$ di offerte del periodo,
dell'ordine del migliaio, il che rende praticabile l'esecuzione su decine di migliaia di
aste.

\begin{lstlisting}[caption={Ricerca del prezzo di equilibrio (\texttt{src/mgp/curve.py}).},
                   label={lst:clearing}]
def prezzo_equilibrio(offerte):
    """Il prezzo di equilibrio e' il piu' piccolo prezzo al quale l'offerta
    cumulata raggiunge la domanda cumulata: p* = min { p : S(p) >= D(p) }."""
    vendita  = offerte[offerte["PURPOSE_CD"] == "OFF"]
    acquisto = offerte[offerte["PURPOSE_CD"] == "BID"]
    if vendita.empty or acquisto.empty:
        return Equilibrio(None, None, None, None, "curva_vuota")

    # Le curve sono costanti a tratti: cambiano valore solo su un prezzo
    # offerto, quindi l'intersezione cade sempre su uno di questi punti.
    prezzi = np.unique(np.concatenate([p_off, p_dom]))
    S = offerta_cumulata_fino_a(prezzi)     # non decrescente
    D = domanda_cumulata_oltre(prezzi)      # non crescente

    incrocio = np.flatnonzero(S >= D)
    if incrocio.size == 0:
        return Equilibrio(None, None, None, None, "domanda_sempre_superiore")

    i = incrocio[0]
    # L'offerta marginale entra solo per la parte che pareggia la domanda.
    return Equilibrio(prezzi[i], min(S[i], D[i]), S[i], D[i], "ok")
\end{lstlisting}

La funzione \`e indipendente dalla granularit\`a: riceve un insieme di offerte gi\`a filtrato
per zona e periodo e non ha bisogno di sapere se il periodo dura quindici, trenta o sessanta
minuti. Questa scelta permette di usare lo stesso codice sull'intero campione, che a cavallo
del 1\textsuperscript{o} ottobre 2025 cambia risoluzione temporale.

\subsection{Il perimetro zonale: la zona NORD come mercato isolato}
\label{sec:perimetro}

L'analisi si concentra sulla zona NORD, modellata come mercato a s\'e stante: le curve sono
costruite con le sole offerte di quella zona, ignorando i limiti di transito con le zone
confinanti. La giustificazione \`e che l'oggetto di studio \`e l'effetto \emph{marginale}
dell'accumulo, cio\`e una differenza fra due scenari calcolati sulla stessa curva, e buona
parte dell'errore di livello si compensa nella differenza.

Questa semplificazione ha per\`o un costo che i dati rendono evidente. La
Tabella~\ref{tab:perimetro} riporta l'errore della ricostruzione al crescere del perimetro
considerato, sul giorno pilota del 31 marzo 2026.

\begin{table}[H]
\centering
\caption{Effetto del perimetro zonale sull'accuratezza della ricostruzione. Giorno pilota
         31 marzo 2026, zona NORD, 96 aste a quarto d'ora.}
\label{tab:perimetro}
\begin{tabular}{@{}lrrr@{}}
\toprule
\textbf{Perimetro} & \textbf{Errore mediano} & \textbf{Scarto mediano} & \textbf{Match $\pm 5$~\eur} \\
                   & (\euromwh{})            & (\euromwh{})            & \\
\midrule
NORD isolata                       & 128{,}43 & $+128{,}43$ & 0\,\% \\
NORD + frontiere confinanti        & 100{,}79 & $+100{,}79$ & 0\,\% \\
\quad + offerte all'altra granularità &  24{,}57 & $+24{,}57$  & 11{,}5\,\% \\
\quad + blocco di scambio netto    & \textbf{5{,}25} & $+1{,}52$ & 50{,}0\,\% \\
\bottomrule
\end{tabular}
\end{table}

Il prezzo ricostruito sulla sola zona NORD risulta \textbf{sistematicamente pi\`u alto} di
quello ufficiale. La ragione \`e strutturale: la zona \`e importatrice, e l'energia che vi
arriva dalle altre zone e dall'estero non compare fra le offerte con codice di zona NORD.
Nel periodo 40 del giorno pilota, per esempio, gli acquisti assegnati in NORD ammontano a
$21\,304$~MW contro $10\,920$~MW di vendite assegnate: oltre 10~GW di differenza entrano
nella zona dall'esterno. Escludendoli si toglie dalla curva di offerta proprio la parte a
basso costo, e il prezzo di equilibrio sale.

\subsection{Il blocco di scambio netto}
\label{sec:scambio}

Per colmare questa lacuna si introduce nelle curve un blocco \emph{price taker} pari alla
differenza fra acquisti e vendite assegnate nel perimetro. Il trattamento \`e simmetrico:
quando la zona importa, il blocco entra come offerta di vendita al prezzo minimo — energia
gi\`a allocata altrove, che si colloca comunque e si comporta quindi da offerta anelastica;
quando la zona esporta, entra come offerta di acquisto al prezzo massimo, cio\`e come domanda
proveniente dall'esterno.

La simmetria non \`e un dettaglio implementativo. In una prima versione era gestito il solo
import, e i dieci periodi con l'errore pi\`u grande del mese di gennaio 2025 risultavano
\emph{tutti} periodi di export, con scarti fino a $-85$~\euromwh{}. Trattando anche l'export
la deviazione standard dello scarto sul mese \`e scesa da 5{,}65 a 1{,}65~\euromwh{}.

\paragraph{Limite da dichiarare.} Lo scambio netto \`e calcolato dall'esito osservato
dell'asta: \`e quindi una grandezza \emph{calibrata}, non prevista dal modello. Il prezzo
resta un esito del modello — emerge dall'incrocio delle curve — ma la simulazione
dell'accumulo assume implicitamente che i flussi con l'esterno non reagiscano alla variazione
di prezzo che l'accumulo induce. \`E un'assunzione forte, tanto pi\`u forte quanto pi\`u il
periodo \`e congestionato, e va richiamata fra i limiti del modello.

Una conseguenza inattesa di questa impostazione \`e che i periodi congestionati non si
ricostruiscono peggio degli altri: su gennaio 2025 l'errore mediano \`e 0{,}17~\euromwh{} nei
611 periodi non congestionati e 0{,}00 nei 133 congestionati, con correlazione nulla fra
ampiezza dello \emph{spread} zonale ed errore. Non si tratta di robustezza del modello: lo
scambio netto \`e calibrato su quantit\`a che sono gi\`a l'esito dell'asta vincolata, quindi
il vincolo di transito entra nel modello sotto forma di quantit\`a osservata anzich\'e di
vincolo esplicito.

\subsection{Le offerte a blocchi}
\label{sec:blocchi}

Le offerte a blocchi sono accettabili solo per intero e su tutti i periodi che coprono, e
trattarle come divisibili introduce un errore che nel mercato a quarto d'ora \`e la principale
fonte di imprecisione: i blocchi sono il 2{,}82\% delle offerte ma valgono in media 385~MW
l'uno contro i 40~MW delle offerte semplici, e spostano l'offerta di circa 880~MW per asta,
il 5\% del volume scambiato.

Per rappresentarne il vincolo \`e stato implementato un clearing iterativo: si parte con tutti
i blocchi accettati, si risolvono tutte le aste della giornata, si rivaluta ciascun blocco sul
prezzo medio ponderato dei periodi che copre — un blocco di vendita resta accettato se il suo
prezzo non supera quel valore — e si ripete fino a stabilit\`a. \`E l'euristica classica per i
mercati con offerte a blocchi; non garantisce l'ottimo, perch\'e il problema esatto \`e un
programma intero, e pu\`o oscillare fra configurazioni.

% DA COMPLETARE: i risultati dell'euristica sono riportati nella Sezione 4.4; resta da
% migliorarne il comportamento nelle giornate in cui oscilla.

\section{Dataset}
\label{sec:dataset}

\subsection{I file delle offerte pubbliche}

I dati sono i file \emph{Offerte Pubbliche} del Mercato del Giorno Prima pubblicati dal GME:
un file XML per ogni giorno di mercato, contenente tutte le zone e tutti i periodi. L'archivio
utilizzato copre il periodo dal 13 febbraio 2015 al 31 marzo 2026, per un totale di 4\,065
giornate; l'anno 2025 \`e completo.

I file recenti sono di dimensioni considerevoli: quello del 31 marzo 2026 pesa 574~MB e
contiene $568\,185$ offerte, di cui $137\,039$ relative alla zona NORD. Un caricamento che
costruisse in memoria l'intero albero XML richiederebbe diversi gigabyte di memoria e un
riparsing completo a ogni esecuzione; la lettura avviene quindi in \emph{streaming}, con
filtro sulla zona applicato durante il parsing e cache del risultato in formato colonnare.

\subsection{Campi rilevanti}

Ogni offerta \`e una coppia prezzo--quantit\`a riferita a un operatore, una zona e un periodo.
I campi utilizzati sono riassunti nella Tabella~\ref{tab:campi}.

\begin{table}[H]
\centering
\caption{Campi dei file di offerte pubbliche utilizzati nell'analisi.}
\label{tab:campi}
\begin{tabular}{@{}lp{8.5cm}@{}}
\toprule
\textbf{Campo} & \textbf{Contenuto} \\
\midrule
\texttt{PURPOSE\_CD}   & Finalità: \texttt{BID} acquisto, \texttt{OFF} vendita \\
\texttt{STATUS\_CD}    & Esito dell'offerta (sei valori, si veda oltre) \\
\texttt{ZONE\_CD}      & Zona di mercato \\
\texttt{PERIOD}        & Periodo del giorno, da leggere insieme a \texttt{GRANULARITY} \\
\texttt{GRANULARITY}   & Granularità: \texttt{PT15}, \texttt{PT30}, \texttt{PT60} \\
\texttt{ENERGY\_PRICE\_NO}    & Prezzo offerto, \euromwh{} \\
\texttt{QUANTITY\_NO}         & Quantità offerta, MW \\
\texttt{AWARDED\_QUANTITY\_NO}& Quantità assegnata, MW \\
\texttt{AWARDED\_PRICE\_NO}   & Prezzo di assegnazione, \euromwh{} \\
\texttt{OFFER\_TYPE}          & \texttt{S} offerta semplice, \texttt{B} offerta a blocchi \\
\texttt{BLOCK\_ID}            & Identificativo del blocco \\
\bottomrule
\end{tabular}
\end{table}

\subsection{Caratteristiche del dato verificate empiricamente}
\label{sec:caratteristiche}

L'ispezione diretta dei file ha corretto diverse assunzioni iniziali. Le si riporta perch\'e
ciascuna ha conseguenze sul trattamento dei dati.

\paragraph{Separatore decimale.} Nei file XML il separatore decimale \`e il \textbf{punto}
(\texttt{146.800}, \texttt{1100.00}), non la virgola: quest'ultima compare negli \emph{export}
in formato Excel o CSV degli stessi dati. La conversione implementata resta comunque
tollerante verso entrambi i separatori, perch\'e l'archivio copre undici anni e una
conversione rigida produrrebbe valori mancanti silenziosi, cio\`e offerte che scompaiono dalle
curve senza segnalazione.

\paragraph{Unit\`a di misura delle quantit\`a.} I campi di quantit\`a sono espressi in
\textbf{MW}, cio\`e sono potenze, e non rappresentano l'energia del periodo. La verifica \`e
indiretta ma conclusiva: sommando le quantit\`a assegnate a livello nazionale in un quarto
d'ora si otterrebbero 148~GW, circa quattro volte il fabbisogno italiano. Il confronto fra un
giorno orario e uno a quarto d'ora conferma la lettura: $38\,334$~MW medi all'ora il
15 gennaio 2025 contro $31\,956$~MW medi al quarto d'ora il 31 marzo 2026, con un rapporto di
0{,}83 e non di 0{,}25. La conseguenza pratica \`e che un'offerta oraria vale la stessa
potenza in ciascuno dei quattro quarti d'ora che compongono l'ora, senza divisioni.

\paragraph{Stato delle offerte.} Il campo di stato assume sei valori e non due:
\texttt{ACC} accettata, \texttt{REJ} rifiutata, \texttt{PREJ} paradossalmente rifiutata,
\texttt{INC} incongrua, \texttt{REP} sostituita da una presentazione successiva,
\texttt{REV} revocata. Le curve d'asta sono costruite con le offerte che hanno effettivamente
partecipato all'asta, cio\`e \texttt{ACC}, \texttt{REJ} e \texttt{PREJ}: includere le
sostituite significherebbe contare due volte la stessa offerta. La scelta \`e confermata dai
dati — includendo anche le sostituite l'errore mediano sul giorno pilota passa da
100{,}79 a 358{,}53~\euromwh{} — e nella Tabella~\ref{tab:stati} sono riportate le varianti
messe a confronto.

\begin{table}[H]
\centering
\caption{Effetto della selezione degli stati sull'accuratezza. Giorno pilota, perimetro NORD
         più frontiere confinanti.}
\label{tab:stati}
\begin{tabular}{@{}lrr@{}}
\toprule
\textbf{Offerte incluse} & \textbf{Aste ricostruite} & \textbf{Errore mediano (\euromwh{})} \\
\midrule
\texttt{ACC} + \texttt{REJ} + \texttt{PREJ} (adottato) & 96 & 100{,}79 \\
\texttt{ACC} + \texttt{REJ}                            & 96 & 100{,}79 \\
solo \texttt{ACC}                                      &  2 & --- \\
\texttt{ACC} + \texttt{REJ} + \texttt{PREJ} + \texttt{REP} & 96 & 358{,}53 \\
tutti gli stati                                        & 96 & 281{,}84 \\
\bottomrule
\end{tabular}
\end{table}

Con le sole offerte accettate le curve non si intersecano quasi mai, ed \`e atteso: le
offerte accettate sono per costruzione quelle interne al mercato, quindi la curva risulta
troncata proprio dove servirebbe il punto di incontro.

\paragraph{Granularit\`a mista.} La granularit\`a non \`e uniforme nemmeno all'interno di una
singola giornata. Nel giorno pilota, zona NORD, convivono $131\,277$ righe a quarto d'ora,
$5\,466$ orarie e 296 semiorarie. Il numero di periodo non \`e quindi interpretabile da solo:
va sempre letto insieme alla granularit\`a, perch\'e con \texttt{PT15} vale da 1 a 96, con
\texttt{PT30} da 1 a 48 e con \texttt{PT60} da 1 a 24. Filtrare sul solo numero di periodo
mescolerebbe quarti d'ora e ore.

\paragraph{Cambio del periodo di mercato.} Il passaggio dall'ora al quarto d'ora cade il
1\textsuperscript{o} ottobre 2025: il 30 settembre la totalit\`a delle offerte \`e oraria, il
giorno successivo l'82{,}8\% \`e a quarto d'ora. L'anno solare 2025 \`e quindi composto da
nove mesi a ventiquattro aste giornaliere e tre mesi a novantasei.

\paragraph{Il prezzo ufficiale \`e nei dati.} Poich\'e il mercato \`e un'asta a prezzo
uniforme, sulle righe accettate il campo del prezzo di assegnazione \`e costante entro
ciascuna coppia zona--periodo e coincide con il prezzo zonale pubblicato. La verifica \`e
positiva su tutti i periodi esaminati. Questo rende disponibile il termine di confronto per la
validazione senza dover allineare dataset aggiuntivi.

\section{Definizione degli Scenari}
\label{sec:scenari}

Uno scenario \`e individuato da tre elementi: la capacit\`a aggregata installata $K$, la
durata dell'accumulo, e il campione di giorni su cui la si simula.

\subsection{La griglia di capacit\`a}

La capacit\`a aggregata varia su una griglia da 1~MW a 6~GW, infittita nella parte bassa.
La scelta degli estremi merita una nota, perch\'e le prime due versioni erano inadeguate.

La prima partiva da 250~MW, ragionando sul volume complessivo del mercato — in zona NORD si
scambiano fra 15 e 25~GW per asta — e ritenendo marginale tutto ci\`o che stava sotto. La
simulazione ha smentito il ragionamento: gi\`a a 250~MW l'erosione risultava di alcuni punti
percentuali e la soglia cadeva fuori griglia. Ci\`o che determina l'effetto sul prezzo non \`e
infatti il volume complessivo ma la \textbf{pendenza della curva attorno all'equilibrio}, dove
i gradini sono pochi e larghi: un accumulo piccolo rispetto al mercato pu\`o essere grande
rispetto alla porzione di curva che effettivamente attraversa.

La seconda partiva da 25~MW e risolveva la soglia principale, ma non quelle definite su
livelli di erosione pi\`u stringenti, che continuavano a cadere sul primo punto. La griglia
definitiva scende fino a 1~MW, il che ha anche permesso di individuare e correggere il
pavimento di discretezza discusso nella Sezione~\ref{sec:pavimento}.

\subsection{La durata}
\label{sec:durata}

Il rapporto fra capacit\`a energetica e potenza determina per quanti periodi consecutivi
l'accumulo pu\`o caricare o scaricare a potenza piena. Il valore di riferimento \`e
\textbf{quattro ore}, il taglio pi\`u diffuso negli impianti in costruzione; durate diverse
sono trattate nell'analisi di sensitivit\`a (Sezione~\ref{sec:sensitivita}).

La durata non \`e un dettaglio tecnico: a parit\`a di potenza, una durata maggiore obbliga
l'accumulo a operare anche in periodi meno favorevoli, spalmando l'effetto su pi\`u ore e
quindi riducendo l'impatto su ciascuna.

La scelta di quattro ore trova un riscontro indipendente in
\textcite{mcconnell2015storage}, che valutano l'accumulo sul mercato all'ingrosso
australiano e trovano \textbf{valore marginale modesto oltre le sei ore} di durata: al di
l\`a di quel punto l'energia aggiuntiva finisce impiegata in periodi troppo poco
remunerativi per giustificarne il costo. Quattro ore stanno dunque nella regione in cui il
rendimento della durata \`e ancora crescente, e le durate maggiori esplorate nell'analisi
di sensitivit\`a servono anche a verificare se la saturazione si collochi, su questo
mercato, dove quel lavoro la trova.

\subsection{I parametri tecnici di riferimento}
\label{sec:parametri}

I parametri fisici dell'accumulo sono ripresi da \textcite{alonsoperez2026storage}, in modo che
non siano scelte arbitrarie di questo lavoro ma valori documentati in letteratura
(Tabella~\ref{tab:parametri}).

\begin{table}[H]
\centering
\caption{Parametri tecnici di riferimento dell'accumulo.}
\label{tab:parametri}
\begin{tabular}{@{}lrl@{}}
\toprule
\textbf{Parametro} & \textbf{Valore} & \textbf{Ruolo nel modello} \\
\midrule
Rendimento di ciclo completo & 92\,\%          & bilancio energetico \\
Costo variabile              & 12~\euromwh{}   & penalizza l'energia scaricata \\
Durata                       & 4 ore           & rapporto capacit\`a/potenza \\
Cicli equivalenti annui      & 350             & solo per il calcolo economico annuo \\
\bottomrule
\end{tabular}
\end{table}

Tre precisazioni evitano altrettanti fraintendimenti.

La condizione «potenza del convertitore pari al 25\% della capacit\`a energetica», che gli
autori indicano come soglia oltre la quale il dimensionamento del convertitore cessa di essere
vincolante, \textbf{coincide} con una durata di quattro ore: $P = 0{,}25\,E$ se e solo se
$E/P = 4$~h. Non \`e quindi un parametro aggiuntivo, ma una conferma indipendente della scelta
gi\`a adottata.

Il \textbf{costo variabile} rappresenta il costo opportunit\`a del degrado, cio\`e la quota di
vita utile che ogni megawattora ciclato consuma. Si applica alla sola scarica, cos\`i che un
ciclo completo lo sconti una volta e non due. La sua introduzione non \`e neutrale: rende
l'accumulo pi\`u selettivo, gli fa compiere meno cicli e ne riduce di conseguenza l'effetto sul
prezzo. Sul campione di gennaio 2025 questo alza la soglia $K^*$ da 73 a 136~MW rispetto a un
modello che ignori il degrado — un modello privo di costo variabile sovrastima l'aggressivit\`a
dell'accumulo e sottostima la capacit\`a che il mercato pu\`o assorbire.

I \textbf{350 cicli annui} non entrano nel programma lineare, che \`e giornaliero: sarebbero un
vincolo annuale imposto a un modello che non vede l'anno. Servono al calcolo economico del
Capitolo~\ref{cap:valutazione}. Corroborano per\`o il vincolo di ciclo chiuso giornaliero, che
ne implica 365. Va segnalato che il modello, ottimizzando in previsione perfetta, ne produce
circa 450: il divario \`e coerente con il fatto che la previsione perfetta fornisce un limite
superiore, e non un risultato conseguibile.

\paragraph{Una previsione sulla sensitivit\`a.} Fra i parametri della
Tabella~\ref{tab:parametri} il rendimento di ciclo \`e quello per cui ci si attende la
minore sensibilit\`a del risultato. \textcite{mcconnell2015storage} osservano che dove il
valore dell'accumulo si concentra nei prezzi estremi, il rendimento round-trip incide
poco: quando il differenziale fra l'ora di carica e quella di scarica \`e ampio, qualche
punto percentuale di perdita non cambia la convenienza dell'operazione, e sposta il piano
solo nelle giornate a margine sottile. Il trasferimento richiede cautela, perch\'e il loro
\`e un mercato \emph{energy only} mentre in Italia esistono un mercato della capacit\`a e
il \ac{MACSE}; l'osservazione vale per\`o per la sola componente di arbitraggio, che \`e
quella considerata qui. Ne discende una previsione verificabile: la sensitivit\`a sul
rendimento (85, 90, 92\,\%) dovrebbe spostare $K^*$ meno di quanto la sposti la durata.

\subsection{Il campione dei giorni}

Il campione \`e costituito da giorni di mercato effettivamente osservati, ciascuno con le sue
curve d'asta. I risultati vengono prodotti sia sul campione complessivo sia
\textbf{stratificati} per stagione e per anno, perch\'e la soglia non \`e stazionaria: dipende
dal mix produttivo e dal livello dei prezzi, che cambiano nel corso dell'anno e fra anni
successivi. Stratificare serve a misurare quanto la soglia si muove, invece di mediare su
condizioni disomogenee.

% DA COMPLETARE: fissare l'estensione definitiva del campione una volta deciso il perimetro
% temporale (si veda la questione aperta sulla non omogeneita' dell'anno 2025).

\section{Esecuzione della simulazione}
\label{sec:esecuzione}

\subsection{Validazione della ricostruzione}

Prima di introdurre l'accumulo occorre stabilire quanto la ricostruzione sia affidabile. Il
criterio adottato \`e la \emph{frequenza di match}: la quota di aste in cui il prezzo
ricostruito coincide con quello ufficiale entro una tolleranza fissata, affiancata
dall'errore assoluto mediano e dal suo segno.

La validazione \`e stata condotta su un mese intero — gennaio 2025, 744 aste orarie — perch\'e
una giornata singola dice se la procedura funziona, non se \`e affidabile. I risultati sono
riportati nella Tabella~\ref{tab:validazione}.

\begin{table}[H]
\centering
\caption{Validazione della ricostruzione su gennaio 2025, zona NORD, 744 aste orarie,
         configurazione adottata.}
\label{tab:validazione}
\begin{tabular}{@{}lr@{}}
\toprule
\textbf{Indicatore} & \textbf{Valore} \\
\midrule
Aste con equilibrio esistente            & 744 su 744 \\
Prezzo ufficiale medio                   & 143{,}32~\euromwh{} \\
Prezzo ricostruito medio                 & 143{,}53~\euromwh{} \\
Mediana dello scarto                     & 0{,}00~\euromwh{} \\
Media dello scarto                       & $+0{,}22$~\euromwh{} \\
Deviazione standard dello scarto         & 1{,}06~\euromwh{} \\
Scarto massimo in valore assoluto        & 6{,}61~\euromwh{} \\
Match entro $\pm 0{,}01$~\euromwh{}      & 52{,}3\,\% \\
Match entro $\pm 1$~\euromwh{}           & 83{,}7\,\% \\
Match entro $\pm 5$~\euromwh{}           & 99{,}5\,\% \\
\bottomrule
\end{tabular}
\end{table}

\`E utile collocare questi valori rispetto alla letteratura.
\textcite{veenstra2025profitability}, ricostruendo il prezzo del mercato olandese del giorno
prima con un impianto analogo, riportano per il 2023 una coincidenza \textbf{esatta nel 32\%
delle ore} e uno scarto \textbf{entro 1~\euromwh{} nell'88\%}. Il confronto con i valori della
Tabella~\ref{tab:validazione} — 52,3\% e 83,7\% — va letto con cautela, perch\'e riguarda
mercati, anni e perimetri diversi: la ricostruzione presentata qui coglie il prezzo esatto
\textbf{pi\`u spesso}, ma \`e \textbf{leggermente meno accurata} nella fascia entro un euro. I
due ordini di grandezza sono comunque comparabili, il che colloca l'accuratezza ottenuta nella
norma di quanto la letteratura consegue con questo tipo di ricostruzione.

% DA COMPLETARE: rifare il confronto sul campione definitivo, e verificare se il quadro cambia
% includendo il regime a quarto d'ora, dove l'accuratezza e' sensibilmente inferiore.

La configurazione a cui si riferiscono questi valori comprende, oltre a quanto descritto
finora, due elementi introdotti dopo una diagnosi degli scarti residui: l'uso della
quantit\`a rettificata al posto di quella offerta, e il trattamento delle offerte a blocchi
come indivisibili. Il contributo di ciascun elemento \`e riportato nella
Tabella~\ref{tab:contributi}.

\begin{table}[H]
\centering
\caption{Contributo dei singoli elementi del modello. Gennaio 2025, 744 aste orarie.}
\label{tab:contributi}
\begin{tabular}{@{}lrrr@{}}
\toprule
\textbf{Configurazione} & \textbf{Dev.\ standard} & \textbf{Match esatto} & \textbf{Match $\pm 1$~\eur} \\
                        & (\euromwh{})            &                       & \\
\midrule
Base (quantit\`a offerta, blocchi divisibili) & 1{,}65 & 44{,}5\,\% & 76{,}7\,\% \\
\quad + quantit\`a rettificata                & 1{,}43 & 49{,}7\,\% & 80{,}8\,\% \\
\quad + blocchi indivisibili                  & \textbf{1{,}06} & \textbf{52{,}3\,\%} & \textbf{83{,}7\,\%} \\
\bottomrule
\end{tabular}
\end{table}

L'accuratezza \`e stabile nel tempo — l'errore mediano giornaliero ha mediana 0{,}11~\euromwh{}
e non supera mai 1{,}81 — ed \`e uniforme nell'arco della giornata. Quest'ultimo punto \`e
rilevante per il seguito: nelle ore di picco serale, quelle in cui un accumulo scaricherebbe,
l'errore mediano resta compreso fra 0{,}00 e 0{,}21~\euromwh{}, in linea con le ore notturne.
Un modello impreciso proprio nelle ore di picco distorcerebbe in modo sistematico la stima del
ricavo da arbitraggio.

\subsection{Il divario fra risoluzioni temporali}

La ricostruzione delle aste a quarto d'ora \`e nettamente meno accurata di quella delle aste
orarie. Il confronto \`e stato condotto sullo stesso mese dell'anno per non confondere
l'effetto della risoluzione con quello della stagionalit\`a.

\begin{table}[H]
\centering
\caption{Confronto fra risoluzioni temporali, zona NORD, configurazione adottata.}
\label{tab:risoluzioni}
\begin{tabular}{@{}lrr@{}}
\toprule
 & \textbf{Gennaio 2025} & \textbf{12--18 gennaio 2026} \\
 & (744 aste orarie)     & (672 aste a quarto d'ora) \\
\midrule
Mediana dello scarto (\euromwh{})        & 0{,}00 & 2{,}52 \\
Deviazione standard (\euromwh{})         & 1{,}06 & 5{,}40 \\
Match esatto ($\pm 0{,}01$~\eur)         & 52{,}3\,\% & 6{,}1\,\% \\
Match entro $\pm 1$~\euromwh{}           & 83{,}7\,\% & 28{,}3\,\% \\
Match entro $\pm 5$~\euromwh{}           & 99{,}5\,\% & 69{,}8\,\% \\
\bottomrule
\end{tabular}
\end{table}

La curva di domanda \`e ricostruita quasi esattamente in entrambi i mercati — lo scarto
rispetto alle quantit\`a assegnate \`e sotto il decimo di punto percentuale — mentre l'errore
sulla curva di offerta \`e circa undici volte pi\`u grande nelle aste a quarto d'ora. La causa
principale sono le offerte a blocchi, che in quel mercato sono cinque volte pi\`u frequenti e
quasi dieci volte pi\`u grandi della media.

\subsection{Attribuzione dello scarto residuo}

Lo scarto fra offerta ricostruita e offerta assegnata si scompone, riga per riga, in tre
contributi la cui somma lo riproduce esattamente: offerte in merito mai assegnate, offerte in
merito assegnate solo in parte, offerte fuori merito comunque assegnate. Attribuendo i
megawatt incoerenti alle loro cause si ottiene la Tabella~\ref{tab:residuo}.

\begin{table}[H]
\centering
\caption{Attribuzione dei megawatt incoerenti con il prezzo del proprio periodo.
         Aste a quarto d'ora, su 990~MW per asta di scarto complessivo.}
\label{tab:residuo}
\begin{tabular}{@{}lrr@{}}
\toprule
\textbf{Causa} & \textbf{MW per asta} & \textbf{Quota} \\
\midrule
Offerte a blocchi trattate come divisibili        & 875{,}8 & 80{,}7\,\% \\
Zone di frontiera senza vincoli di transito       &  76{,}6 &  7{,}1\,\% \\
Offerte con quota minima di accettazione          &   0{,}0 &  0{,}0\,\% \\
Non spiegato                                      & 132{,}8 & 12{,}2\,\% \\
\bottomrule
\end{tabular}
\end{table}

Le prime due cause sono semplificazioni dichiarate del modello, non errori, e il loro costo
\`e ora misurato. La terza merita una nota: il campo che indica la quota minima di
accettazione risulta valorizzato su tutte e sole le offerte a blocchi — la tabella di
contingenza \`e perfettamente diagonale — quindi non \`e un meccanismo separato ma un
attributo dei blocchi, gi\`a conteggiato nella prima riga.

Restano non spiegati 133~MW per asta, pari allo 0{,}74\% del volume scambiato. Due ipotesi
sono state formulate e scartate sulla base dei dati: l'accettazione parziale delle offerte,
che vale lo 0{,}25\% dei megawatt assegnati con frequenza simile nei due mercati, e
l'esistenza di vincoli che legano quarti d'ora consecutivi. Su quest'ultima il test
discriminante riguarda le offerte identiche nei quattro quarti il cui prezzo cade fra il
minimo e il massimo dei quattro prezzi dell'ora, cio\`e le uniche per cui la decisione
corretta cambierebbe da un quarto all'altro: il 99{,}3\% di esse viene assegnato in modo
diverso fra i quarti, come prevede un clearing indipendente.

% DA COMPLETARE: il 12,2% non spiegato resta da attribuire.

\subsection{Il calcolo dell'erosione, giorno per giorno}
\label{sec:calcolo-erosione}

Per ciascun giorno del campione e ciascuna capacit\`a della griglia la procedura \`e quella
definita nella Sezione~\ref{sec:erosione}:

\begin{enumerate}
  \item si ricostruiscono le curve d'asta di tutti i periodi della giornata e se ne ricava il
        prezzo di riferimento $p^{0}_t$, cio\`e quello senza accumulo;
  \item si ottimizza \emph{una volta sola} il piano di carica e scarica su quei prezzi;
  \item lo si valorizza ai prezzi di riferimento, ottenendo $\pi_{PT}$;
  \item si inserisce lo stesso piano nelle curve, si ricalcola l'equilibrio di ogni periodo e
        lo si valorizza ai prezzi nuovi, ottenendo $\pi_{PM}$;
  \item la differenza fornisce l'erosione, in valore assoluto e in quota.
\end{enumerate}

Le curve e i prezzi di riferimento si calcolano una sola volta per giornata e si riusano per
tutte le capacit\`a: cambia soltanto il profilo che vi si inserisce.

\subsection{Ispezione delle curve nei periodi estremi}
\label{sec:ispezione}

Il risultato dipende interamente dalla pendenza della curva di offerta nelle ore in cui
l'accumulo opera, cio\`e quella di minimo e quella di massimo prezzo. Poich\'e una pendenza
sostenuta da poche offerte isolate renderebbe il risultato fragile, le curve sono state
ispezionate direttamente (Figura~\ref{fig:curve}).

\begin{figure}[H]
\centering
\includegraphics[width=\textwidth]{07_curve_20250120}
\caption{Curve aggregate di domanda e offerta ricostruite per l'ora di minimo (in alto) e
         quella di picco (in basso) del 20 gennaio 2025, in vista completa e ingrandite
         attorno all'equilibrio. I punti sulla curva di offerta sono i singoli gradini.}
\label{fig:curve}
\end{figure}

L'ispezione conferma tre cose. La pendenza attorno all'equilibrio poggia su \textbf{decine di
gradini}: nella finestra ingrandita ne cadono fra 39 e 62, ed entro $\pm 20$~\euromwh{} dal
prezzo di equilibrio si contano fra 42 e 72 offerte di vendita per 3.500-5.500~MW, la maggiore
delle quali vale circa un decimo del totale locale. Le offerte presentate ai \textbf{limiti di
prezzo} — il blocco a $-500$~\euromwh{} che occupa i primi migliaia di megawatt e la domanda
price taker a 4.000 — stanno agli estremi della curva e non toccano la regione di equilibrio.
Il prezzo ricostruito \textbf{cade dove atteso}: coincide con quello ufficiale al centesimo in
tre dei quattro periodi ispezionati, e nel quarto differisce di 1{,}50~\euromwh{} in un
periodo in cui la curva di offerta \`e visibilmente pi\`u piatta attorno all'incrocio.

Una caratteristica emerge con chiarezza dai grafici: la \textbf{curva di domanda \`e quasi
verticale} in corrispondenza dell'equilibrio, in tutti i periodi ispezionati. L'effetto
dell'accumulo passa quindi quasi interamente attraverso la curva di offerta, e ci\`o spiega
perch\'e sia la pendenza di quest'ultima a determinare la soglia.

\subsection{La curva di impatto marginale}
\label{sec:impatto}

Accanto al ricalcolo esatto dell'equilibrio, con cui si producono i risultati, \`e utile una
rappresentazione che renda leggibile in un colpo d'occhio quanto un periodo sia elastico: la
curva $\Delta p = f(\Delta q)$ della variazione di prezzo in funzione della quantit\`a aggiunta
o sottratta, con l'\textbf{origine centrata sul punto di clearing osservato} del periodo. La
tecnica \`e ripresa da \textcite{alonsoperez2026storage}, che la usano come proxy
dell'elasticit\`a del mercato spagnolo.

Si legge cos\`i: il ramo destro ($\Delta q > 0$) \`e l'accumulo che scarica, aggiunge offerta e
abbassa il prezzo; quello sinistro \`e l'accumulo che carica, aggiunge domanda e lo alza. I
tratti piatti sono i gradini larghi, dove l'accumulo passa inosservato; i salti sono i confini
fra un gradino e il successivo. L'asimmetria fra i due rami \`e la ragione per cui l'effetto
dell'accumulo non si riassume in un unico numero di elasticit\`a.

\begin{figure}[H]
\centering
\includegraphics[width=\textwidth]{07_impatto_20250120}
\caption{Curva di impatto marginale dell'ora di minimo e dell'ora di picco del 20 gennaio
         2025. L'origine \`e il prezzo di equilibrio effettivamente ricostruito per quel
         periodo.}
\label{fig:impatto}
\end{figure}

La Figura~\ref{fig:impatto} rende visibile il meccanismo che governa l'intero risultato.
Attorno all'origine si estende un \textbf{pianoro} largo alcune centinaia di megawatt in cui il
prezzo non si muove affatto: l'equilibrio cade in mezzo a un gradino, e finch\'e l'accumulo non
lo esaurisce il prezzo marginale resta lo stesso. Superato il gradino la curva scende a
scalini, ciascuno dei quali \`e un'offerta che entra o esce dalla soluzione. \`E questa
struttura — non l'ampiezza del mercato — a determinare la soglia: ci\`o che conta \`e il
rapporto fra la taglia dell'accumulo e la larghezza dei gradini che attraversa.

La curva non \`e ottenuta ripetendo il clearing a ogni punto, ma invertendo la funzione di
eccesso di offerta. Aggiungere una quantit\`a $q$ di offerta price taker trasla $S(p)-D(p)$
verso l'alto di $q$ a ogni prezzo, da cui
\begin{equation}
p^{*}(q) \;=\; \min\left\{\, p \;:\; S(p) - D(p) \ge -q \,\right\},
\label{eq:impatto}
\end{equation}
che \`e esatta, non approssimata, e sfrutta la monotonia dell'eccesso di offerta gi\`a
richiamata nella Sezione~\ref{sec:mgp}.

Due limiti vanno dichiarati. La curva \`e una diagnostica \textbf{per periodo}, costruita a
curve fisse: non rivaluta le offerte a blocchi, che sono indivisibili e coprono pi\`u periodi,
quindi per capacit\`a elevate il prezzo che vi si legge pu\`o discostarsi da quello del
clearing di giornata. E l'asse in energia dipende dalla durata del periodo: su aste a quarto
d'ora un megawatt vale $0{,}25$~MWh, e le due scale coincidono soltanto sulle aste orarie.

\subsection{Il pavimento di discretezza}
\label{sec:pavimento}

Le curve ricostruite sono funzioni a gradini, e questo introduce un effetto che va corretto
prima di interpretare i risultati alle capacit\`a pi\`u piccole.

Con una griglia fitta nella parte bassa si osserva che fra 1 e 15~MW l'erosione cresce
\textbf{molto meno che proporzionalmente}: passa dallo 0{,}7\% all'1{,}3\% in mediana, cio\`e
raddoppia appena mentre la capacit\`a si moltiplica per quindici. Un effetto di mercato
dovrebbe invece scalare con la capacit\`a installata, e il valore a 1~MW dovrebbe essere
prossimo a zero.

L'esame periodo per periodo mostra la ragione: a 1~MW il prezzo si muove mediamente in due
periodi e mezzo sui tredici in cui l'accumulo opera, ma quando si muove \textbf{salta di un
gradino intero}, in media 1{,}4~\euromwh{}. Non \`e uno spostamento graduale: sono gli
equilibri che cadono esattamente sul bordo di un gradino e che pochi megawatt fanno scattare a
quello successivo.

Nel mercato reale il fenomeno non si presenta allo stesso modo, perch\'e l'offerta marginale
viene accettata parzialmente e l'incrocio cade all'interno del gradino. Si tratta quindi di un
\textbf{artefatto della ricostruzione}, non di un fenomeno da misurare.

Il pavimento viene misurato alla capacit\`a pi\`u piccola della griglia — dove per costruzione
non pu\`o esserci effetto di mercato — e sottratto giorno per giorno. Vale lo 0{,}7\% in
mediana e il 2{,}8\% al 90\textsuperscript{o} percentile. La correzione sposta la soglia in
modo sistematico, e in misura tanto maggiore quanto pi\`u basso \`e il livello di erosione
scelto: a un livello del 20\% la sposta del 21\%, al 10\% del 75\%, al 5\% del 117\%. Ne segue
che una soglia fissata al 5\% \`e la \textbf{meno affidabile}, perch\'e il pavimento vale pi\`u
della met\`a della soglia stessa.

\subsection{Dal campione alla soglia: il bootstrap}
\label{sec:bootstrap}

La soglia $K^*$ si stima ricampionando i giorni con reimmissione. Per ciascun
ricampionamento si calcola, a ogni capacit\`a della griglia, il quantile prudenziale
dell'erosione fra i giorni estratti; si ottiene cos\`i una curva che lega il quantile alla
capacit\`a, crescente perch\'e pi\`u accumulo erode di pi\`u; e si individua per interpolazione
la capacit\`a alla quale quella curva attraversa il livello di erosione dichiarato. La
distribuzione delle $K^*$ cos\`i ottenute fornisce l'intervallo di confidenza.

Due scelte vanno dichiarate perch\'e non sono stimate dai dati. Il \textbf{livello di
erosione} che definisce il passaggio a price maker \`e convenzionale: si adotta il 10\%, e la
sensibilit\`a a questa scelta \`e riportata nel Capitolo~\ref{cap:valutazione}. Il
\textbf{quantile} \`e l'80\textsuperscript{o} o il 90\textsuperscript{o} percentile, e non uno
pi\`u estremo: con qualche centinaio di giorni un quantile al 95\textsuperscript{o} o al
99\textsuperscript{o} percentile poggerebbe su pochissime osservazioni e la sua stima sarebbe
instabile proprio nella parte che determina il risultato.

\subsection{Le giornate a basso differenziale}
\label{sec:basso-spread}

Quando il differenziale di prezzo giornaliero \`e piccolo, $\pi_{PT}$ tende a zero e il
rapporto~\eqref{eq:erosione} diventa privo di significato o esplode. Scartare quelle giornate
sarebbe per\`o metodologicamente pericoloso: sono le giornate ad alta produzione rinnovabile
e prezzi compressi, cio\`e proprio quelle destinate a diventare pi\`u frequenti, e escluderle
introdurrebbe un \textbf{bias sistematico} nella direzione opposta a quella dell'evoluzione
del sistema.

La scelta \`e quindi di mantenerle nel campione e di riportare, accanto all'erosione
relativa, anche quella in \textbf{valore assoluto} (euro al giorno), che resta interpretabile
a qualunque livello di differenziale. Il rapporto viene semplicemente omesso dove il profitto
di riferimento \`e trascurabile, e il numero di casi in cui ci\`o accade \`e riportato con i
risultati.

\section{La fase previsiva e la struttura del suo errore}
\label{sec:previsione}

La fase 1 dell'impianto a due fasi (Sezione~\ref{sec:due-fasi}) produce, per ogni giorno del
campione, la previsione delle ventiquattro ore successive formulata con la sola informazione
disponibile alla chiusura del mercato. Questa sezione ne descrive la specifica e, soprattutto,
\textbf{caratterizza l'errore che ne risulta}: non per misurare l'accuratezza del modello, che
non \`e l'oggetto della tesi, ma perch\'e \`e la struttura di quell'errore a propagarsi al
piano dell'accumulo e al risultato economico.

\subsection{L'orizzonte, e perch\'e \`e esattamente ventiquattro}

Il \ac{MGP} del giorno $D$ chiude a mezzogiorno di $D-1$, ma gli esiti di $D-1$ sono stati
pubblicati il giorno prima: al momento della decisione l'operatore conosce \textbf{tutte} le
ventiquattro ore di $D-1$ e nessuna di $D$. L'origine della previsione \`e quindi la fine di
$D-1$ e l'orizzonte $h = 1, \dots, 24$. Nessuna informazione dal futuro entra nel piano, e non
serve troncare la giornata a met\`a.

Si prevedono i prezzi zonali \textbf{ufficiali} pubblicati dal \ac{GME}, non quelli
ricostruiti dal modello di clearing. \`E quello che osserva un operatore reale, e prevedere la
propria ricostruzione significherebbe prevedere anche il proprio errore di ricostruzione, che
non \`e un fenomeno di mercato. La scelta non intacca l'erosione, perch\'e la fase 2 continua
a valorizzare sui prezzi ricostruiti con e senza accumulo: l'errore di ricostruzione entra
identico nei due termini e si semplifica.

\subsection{La specifica}
\label{sec:specifica}

Si adotta un modello \ac{SARIMAX} nella formulazione classica di
\textcite{box2015timeseries}, e non un metodo di apprendimento automatico. La ragione discende
dall'oggetto: poich\'e interessa la \emph{struttura} dell'errore e non la sua sola dimensione,
serve un modello con una teoria dell'errore esplicita --- intervalli di previsione derivati,
residui analizzabili --- e non soltanto misurabile a posteriori. \textcite{weron2014previsione}
documenta del resto che i modelli di questa famiglia restano il termine di confronto
obbligato nella previsione dei prezzi elettrici, qualunque sia il metodo che li affianca.

La specifica adottata \`e
\[
  \text{SARIMAX}(2,1,1)\times(0,1,1)_{24}
\]
con regressori esogeni deterministici.

\paragraph{La parte stagionale non \`e una scelta ma una lettura.} Sulla serie oraria della
zona NORD, la \ac{ACF} della sola differenza prima ai multipli di ventiquattro vale
$+0{,}62$, $+0{,}57$, $+0{,}53$, $+0{,}53$, $+0{,}52$: \textbf{non decade}, cio\`e resta una
stagionalit\`a non rimossa. Aggiungendo la differenza stagionale i medesimi valori diventano
$-0{,}43$, $-0{,}01$, $-0{,}05$, $+0{,}00$, $-0{,}06$: una \textbf{sola punta negativa a
ritardo ventiquattro che poi taglia netto}, con \ac{PACF} pari a $-0{,}429$. \`E la firma di
una media mobile stagionale di ordine uno, da cui $D = 1$, $P = 0$, $Q = 1$, $s = 24$.

\paragraph{La stagionalit\`a settimanale entra come regressore, non come seconda
stagionalit\`a.} Un modello SARIMA ne ammette una sola, e differenziare a ritardo
centosessantotto costerebbe una settimana di osservazioni iniziali. Si usano invece due coppie
di termini di Fourier a periodo $168$ e tre indicatrici, per sabato, domenica e festivi
italiani. La scelta \`e sostenuta dai dati: tolto il profilo orario, il giorno della settimana
spiega il \SI{5,2}{\percent} della varianza residua, con un'escursione di
15,8~\euromwh{}. Merita qualche regressore, non una struttura stagionale.

\paragraph{Nessuna trasformazione logaritmica.} Il minimo osservato nel biennio vale
0,10~\euromwh{} e sul \ac{MGP} i prezzi negativi sono ammessi fino a
$-500$~\euromwh{}: il logaritmo sarebbe fragile per costruzione.

\paragraph{Ordine congelato.} Gli ordini non stagionali sono stati scelti confrontando per
\ac{AIC} quattro candidati dichiarati in anticipo, \emph{una volta sola} sull'anno di
addestramento, e poi \textbf{congelati} per tutto il periodo di valutazione. Sono cio\`e un
dato dichiarato del lavoro, non un iperparametro riottimizzato a ogni passo.

\paragraph{Protocollo a origine mobile.} Per ogni giorno del campione si producono
ventiquattro passi con origine alla fine del giorno precedente. I coefficienti si ristimano
una volta al mese, non ogni giorno: sia perch\'e trecentosessantasei stime complete
costerebbero giorni di calcolo per un guadagno previsivo trascurabile, sia perch\'e nessun
operatore ricalibra un modello ogni notte. Fra una ristima e l'altra il filtro assorbe il
giorno appena osservato aggiornando lo \emph{stato} senza toccare i \emph{coefficienti}.

\subsection{La struttura dell'errore, in sei dimensioni}
\label{sec:errore}

Sia $e(g,h)$ l'errore del giorno $g$ all'ora $h$, definito come prezzo reale meno prezzo
previsto: positivo dove il modello ha sottostimato. Sull'anno di valutazione l'errore vale
14,75~\euromwh{} di \ac{RMSE} e 10,54 di \ac{MAE}. Sono numeri che da soli
non dicono nulla di utile: quello che conta \`e come si distribuiscono.

\paragraph{1. Per ora del giorno.} L'errore non \`e omogeneo: varia di \textbf{2,4 volte} fra
l'ora migliore e la peggiore, da 8,22~\euromwh{} a mezzanotte a 19,97 alle
quattordici (Figura~\ref{fig:errore-orario}). Le ore notturne, piatte, si prevedono bene; il
pomeriggio male.

Il fattore che lo spiega non \`e il livello del prezzo ma la sua \textbf{variabilit\`a}: la
correlazione fra l'errore orario e la deviazione standard del prezzo in quell'ora vale
$+0{,}931$, contro $+0{,}185$ con il livello medio. Il picco serale, che \`e l'ora di prezzo
pi\`u alto dell'anno, si prevede \emph{meglio} del ventre pomeridiano, perch\'e \`e alto ma
regolare, mentre il pomeriggio \`e mediamente basso ma erratico, dipendendo dal fotovoltaico.

Il rilievo per l'arbitraggio \`e diretto: la batteria colloca spesso la carica proprio nel
ventre pomeridiano, che \`e l'ora peggio prevista dell'intera giornata.

\paragraph{2. Distorsione contro dispersione.} Il modello \`e \textbf{non distorto}: media
dell'errore $0{,}015$~\euromwh{}, con un test $t$ sulla media nulla che d\`a $p = 0{,}93$. Il
massimo bias per ora del giorno vale $0{,}29$~\euromwh{} contro una dispersione di
$14{,}75$. Tutto l'errore \`e dunque dispersione, non distorsione sistematica: non esiste una
correzione da applicare, si pu\`o solo convivere con la varianza.

\paragraph{3. Code e forma.} La distribuzione \`e simmetrica (asimmetria $-0{,}0002$) ma
decisamente \textbf{leptocurtica}: curtosi in eccesso $+3{,}97$
(Figura~\ref{fig:errore-forma}). Gli eventi oltre tre deviazioni standard sono \textbf{6,2
volte} pi\`u frequenti dell'attesa gaussiana, quelli oltre quattro \textbf{101 volte}.
L'\SI{1}{\percent} peggiore delle ore porta il \SI{19,0}{\percent} della somma dei quadrati,
il \SI{5}{\percent} peggiore il \SI{45,4}{\percent}.

Questo si lega alla redditivit\`a in modo diretto. L'errore quadratico medio giornaliero
correla con il differenziale del giorno a $r = +0{,}600$ e con la volatilit\`a a $r =
+0{,}638$, mentre con il livello medio dei prezzi non correla affatto ($r = +0{,}045$, non
significativo). L'errore \`e cio\`e massimo proprio nelle giornate in cui l'arbitraggio vale
di pi\`u. Il modello inoltre \textbf{contrae il differenziale} verso la sua media: prevede
valori compresi fra 40 e 80~\euromwh{} quando il reale spazia da 20 a
158, e lo sottostima nel \SI{65}{\percent} delle giornate
(Figura~\ref{fig:errore-spread}).

\paragraph{4. Autocorrelazione.} L'\ac{ACF} dei residui vale $+0{,}849$ a ritardo uno,
$+0{,}670$ a due, $+0{,}518$ a tre; la \ac{PACF} vale $+0{,}849$ a uno, $-0{,}184$ a due e
poi si annulla. Il test di Ljung-Box rifiuta l'ipotesi di residui incorrelati con $p$
inferiore a $10^{-4}$ a tutti i ritardi esaminati. Il trattamento di questo limite \`e nella
Sezione~\ref{sec:limiti-previsione}.

La memoria vive per\`o \textbf{dentro} la giornata e non fra giornate: l'\ac{ACF} dell'errore
medio giornaliero vale $+0{,}034$ a ritardo uno, dentro la banda di non significativit\`a. Una
giornata difficile non ne annuncia un'altra. \`E il verso peggiore per il piano: un errore
coerente per l'intera giornata sbaglia l'\emph{ordinamento} delle ore, mentre errori
indipendenti si compenserebbero fra loro.

\paragraph{5. Eteroschedasticit\`a.} Rispetto al \emph{livello} del prezzo previsto l'errore
segue un andamento \textbf{a U}: \ac{RMSE} pari a $22{,}01$ nel decile pi\`u basso,
$10{,}96$ nei centrali, $20{,}63$ nel pi\`u alto. Una correlazione di rango non lo
rivelerebbe --- vale $+0{,}001$, non significativa --- perch\'e cerca monotonia dove la
relazione non \`e monotona.

Rispetto alla \textbf{volatilit\`a della giornata} la relazione \`e invece monotona e forte:
l'errore medio passa da $9{,}93$ nel quintile pi\`u tranquillo a $20{,}45$ nel pi\`u agitato.
\`E l'anello che lega l'errore di previsione al regime di mercato, e va nella direzione
sfavorevole: i regimi volatili sono quelli in cui l'accumulo guadagna.

I decili si costruiscono sul prezzo \emph{previsto} e non su quello realizzato. Raggruppare
per il realizzato sarebbe contaminato per costruzione, perch\'e nel gruppo dei prezzi alti
finirebbero per definizione le ore in cui il prezzo \`e stato pi\`u alto del previsto: si
misurerebbe una regressione verso la media, non l'eteroschedasticit\`a.

\paragraph{6. Calibrazione degli intervalli.} La copertura empirica degli intervalli di
previsione, confrontata con quella nominale, decresce in modo monotono e cambia segno:
\SI{64,7}{\percent} contro un nominale del \SI{50}{\percent}, 88,2 contro
80, 93,6 contro 90, 96,0 contro 95,
98,1 contro 99 (Figura~\ref{fig:calibrazione}).

\`E la firma esatta di intervalli gaussiani costruiti su un errore leptocurtico: troppo larghi
al centro, troppo stretti nelle code. Il difetto non \`e di \emph{ampiezza} ma di
\emph{forma}, ed \`e la stessa cosa vista al punto 3, letta dal lato dell'incertezza
dichiarata. La miscalibrazione \`e inoltre strutturata per ora: la copertura al
\SI{90}{\percent} vale 97--\SI{99}{\percent} nelle ore notturne e scende a
86,9--\SI{88,8}{\percent} fra le tredici e le quindici, cio\`e il modello \`e pi\`u
sicuro di quanto dovrebbe proprio nelle ore che contano per l'arbitraggio.

\subsection{Due limiti dichiarati}
\label{sec:limiti-previsione}

\paragraph{I residui non sono bianchi, e va detto per primo.} L'autocorrelazione dei residui
vale $+0{,}849$ a ritardo uno e il test di Ljung-Box rifiuta con $p$ inferiore a $10^{-4}$.
Resta quindi \textbf{struttura non catturata}: il modello \`e migliorabile in senso
statistico, e sarebbe scorretto presentarlo come adeguato.

Questo per\`o \textbf{non intacca il risultato economico}, e la ragione non \`e una
minimizzazione ma una conseguenza di quanto stabilito nella Sezione~\ref{sec:ordinale}.
L'autocorrelazione residua riguarda il \emph{livello} dei prezzi --- \`e su quello che il
modello lascia segnale non sfruttato --- mentre il valore dell'arbitraggio dipende
dall'\emph{ordinamento} delle ore. Un modello con residui pi\`u bianchi produrrebbe previsioni
di livello migliori, ma ridurrebbe il costo del \SI{10,4}{\percent} solo nella misura in cui
migliorasse specificamente la classifica delle ore, che \`e una propriet\`a diversa e che
nessuno dei criteri usuali di adattamento misura.

Ne segue anche un'indicazione per chi volesse proseguire: raffinare il modello sul criterio
d'informazione o sulla bianchezza dei residui \`e la strada meno promettente. Converrebbe
piuttosto valutarlo direttamente sulla correlazione di rango giornaliera, che \`e la
statistica che determina il piano.

\paragraph{Ora del giorno e orizzonte non sono separabili.} L'origine della previsione \`e
sempre la fine di $D-1$, quindi l'orizzonte $h$ corrisponde sempre e solo all'ora $h-1$ del
giorno $D$: le due variabili sono \textbf{collineari per costruzione} e nessuna analisi su
questi dati pu\`o attribuire all'una o all'altro lo schema osservato. Non \`e un difetto del
disegno statistico ma della struttura del mercato, che ammette una sola sessione al giorno.

Si pu\`o per\`o escludere che si tratti di un puro effetto di orizzonte. Un errore crescente
con la distanza temporale sarebbe monotono, e lo schema non lo \`e: vale
8,22~\euromwh{} alla prima ora, sale a 19,97 alla quindicesima e ridiscende
a 10,65 alla ventiquattresima. Il coefficiente di correlazione fra errore orario e
variabilit\`a del prezzo in quell'ora vale $+0{,}931$, contro $+0{,}185$ con il livello: il
fattore dominante \`e l'ora, non la distanza.

\begin{figure}[H]
\centering
\includegraphics[width=0.92\textwidth]{15_errore_orario}
\caption{Errore di previsione per ora del giorno. In alto l'errore realizzato, in due
metriche, contro quello \emph{dichiarato} dal modello: il secondo cresce in modo monotono con
l'orizzonte, come impone la teoria, mentre il primo segue l'ora del giorno. Al centro la
variabilit\`a del prezzo in ciascuna ora, che ha la stessa forma dell'errore ($r = +0{,}931$);
in basso il livello medio, che ne ha una diversa ($r = +0{,}185$).}
\label{fig:errore-orario}
\end{figure}

\begin{figure}[H]
\centering
\includegraphics[width=\textwidth]{15_errore_forma}
\caption{Forma della distribuzione dell'errore. L'istogramma da solo ingannerebbe, perch\'e le
code sono rare proprio dove contano; il diagramma quantile-quantile le rende leggibili come
scostamento dalla bisettrice.}
\label{fig:errore-forma}
\end{figure}

\begin{figure}[H]
\centering
\includegraphics[width=\textwidth]{15_errore_spread}
\caption{A sinistra, l'errore giornaliero cresce con il differenziale della giornata: il
modello sbaglia di pi\`u dove l'arbitraggio vale di pi\`u. A destra, il differenziale
previsto resta confinato in una fascia stretta mentre quello reale spazia molto di pi\`u:
il modello lo contrae verso la sua media.}
\label{fig:errore-spread}
\end{figure}

\begin{figure}[H]
\centering
\includegraphics[width=0.78\textwidth]{15_calibrazione}
\caption{Copertura empirica contro copertura nominale degli intervalli di previsione. Lo
scarto decresce in modo monotono e cambia segno: \`e la firma di intervalli gaussiani su un
errore a code spesse.}
\label{fig:calibrazione}
\end{figure}

\section{La soglia $K^*$}
\label{sec:soglia-definitiva}

La soglia si stima sulla griglia definitiva a 132 capacit\`a, con l'erosione al netto del
pavimento di discretezza (Sezione~\ref{sec:pavimento}) e le giornate a piano vuoto contate a
erosione nulla. Si riportano \textbf{due varianti}: piano da previsione perfetta, che \`e il
riferimento, e piano da previsione (Sezione~\ref{sec:due-varianti}).

\subsection{La verifica preliminare di monotonia}

La regola con cui si individua $K^*$ prende il \emph{primo} superamento della soglia da parte
della curva quantile, il che presuppone che la curva sia crescente. Con una griglia fitta un
sobbalzo locale del quantile dentro un ricampionamento potrebbe farla scattare in anticipo e
abbassare la stima: la presunzione va quindi verificata prima di usare il risultato, non dopo.

Su mille ricampionamenti e per tutte e quattro le combinazioni di variante e soglia, la quota
di ricampionamenti con \textbf{pi\`u di un attraversamento \`e nulla}, il calo massimo della
curva vale $-0{,}01$ punti percentuali in mediana, e lo scarto fra ultimo e primo
attraversamento \`e \textbf{esattamente zero}, sia in media sia al massimo. La regola del
primo attraversamento resta quindi valida e non serve alcuna correzione.

\subsection{Il risultato}

\begin{table}[H]
\centering
\caption{Soglia $K^*$ al 90\textsuperscript{o} percentile dell'erosione netta, con intervallo
di confidenza al \SI{90}{\percent} da mille ricampionamenti dei giorni.}
\label{tab:soglia}
\begin{tabular}{@{}llrr@{}}
\toprule
\textbf{Piano} & \textbf{Soglia} & \textbf{$K^*$ (MW)} & \textbf{IC \SI{90}{\percent}} \\
\midrule
previsione perfetta  & \SI{10}{\percent} & 75,9  & 70,5 -- 81,7 \\
previsione perfetta  & \SI{20}{\percent} & 173,1 & 157,5 -- 189,0 \\
da previsione        & \SI{10}{\percent} & 50,1  & 43,3 -- 57,5 \\
da previsione        & \SI{20}{\percent} & 129,3 & 108,0 -- 142,9 \\
\bottomrule
\end{tabular}
\end{table}

\textbf{La differenza fra le due varianti \`e il risultato}: la soglia scende del
\SI{34,0}{\percent} al livello del \SI{10}{\percent} e del \SI{25,3}{\percent} al
\SI{20}{\percent}. Gli intervalli di confidenza delle due varianti non si sovrappongono a
nessuna delle due soglie, quindi la differenza non \`e un artefatto campionario.

Una flotta che pianifica su previsioni realistiche diventa price maker \textbf{prima}. Lo
stesso errore che le costa il \SI{10,4}{\percent} di profitto la fa anche operare in modo meno
mirato, quindi con pi\`u effetto sul prezzo a parit\`a di capacit\`a installata --- e lo
conferma anche l'erosione in valore assoluto, non solo il rapporto.

\paragraph{Quanto pesa il pavimento.} Sulla stessa griglia, ma senza sottrarre il pavimento di
discretezza, la soglia della variante realistica risulterebbe di \SI{33,7}{MW} invece di
50,1. La correzione del pavimento non \`e quindi un raffinamento marginale: senza di
essa la soglia sarebbe sottostimata di un terzo.

\paragraph{Una sensibilit\`a da dichiarare.} Il pavimento \`e mal definito nelle giornate in
cui il profitto di riferimento \`e quasi nullo, dove il rapporto \`e dominato dal rumore: due
giornate su trecentosessantasei superano il \SI{25}{\percent} di pavimento, e in una di esse
il profitto a \SI{1}{MW} vale meno di cinque euro. Escludendo le giornate a pavimento anomalo
il divario fra le due varianti resta compreso fra il 26 e il \SI{34}{\percent}, quindi
il risultato qualitativo \`e robusto, ma il valore puntuale ne dipende. Si adotta la stima sul
campione completo, coerentemente con la scelta di non scartare le giornate a basso
differenziale (Sezione~\ref{sec:basso-spread}), dichiarando la sensibilit\`a.

\begin{figure}[H]
\centering
\includegraphics[width=0.95\textwidth]{17_curva_erosione}
\caption{Curva erosione-capacit\`a nelle due varianti di piano, mediana e
90\textsuperscript{o} percentile, su scala logaritmica delle capacit\`a. La curva della
variante realistica sta sopra quella della previsione perfetta a ogni capacit\`a: i punti
segnano i quattro $K^*$ della Tabella~\ref{tab:soglia}.}
\label{fig:curva-erosione}
\end{figure}

\subsection{La propagazione dell'errore al piano}
\label{sec:propagazione}

Il confronto fra le due varianti permette di misurare separatamente le due perdite. A
capacit\`a piccola, dove l'effetto sul prezzo \`e trascurabile, la differenza fra i due piani
\`e \textbf{pura perdita informativa}: vale il \SI{10,4}{\percent} del limite superiore,
cio\`e un'efficienza informativa del \SI{89,6}{\percent}.

\`E molto meno di quanto la dimensione dell'errore facesse temere, e la
Sezione~\ref{sec:ordinale} ne d\`a la ragione. Due misure la rendono concreta: il modello
individua l'ora esatta del minimo giornaliero solo nel \SI{24,6}{\percent} dei casi, ma sceglie
correttamente il \SI{78,4}{\percent} delle ore di carica e l'\SI{86,2}{\percent} di quelle di
scarica. E in nessuna delle giornate del campione la previsione porta a rinunciare a operare,
n\'e a operare quando non converrebbe.

\paragraph{Un'illusione che sull'aggregato inganna.} La differenza fra il profitto che la
batteria si attende --- il piano valorizzato ai prezzi previsti --- e quello che realizza vale
in totale meno dell'\SI{1}{\percent}, il che farebbe pensare che l'operatore sappia prevedere
il proprio guadagno. \`E \textbf{compensazione, non accuratezza}: giornata per giornata quella
differenza ha media assoluta relativa del \SI{48,1}{\percent}, con errori nei due versi che si
annullano sull'anno. Riportare il solo totale sarebbe fuorviante.
